\section{Methods}

\subsection{Architecture}
\label{sec:arch}
We propose a method that encodes samples as classifications mapped to
contrastive vectors. This directly constrains a model's representations,
and classifications can natively be used for steering.

The model (Figure~\ref{fig:onto-stack}) consists of $l$ $\dictenc$
layers writing into a shared residual accumulator, decoded by a single
linear map. Each layer (Figure~\ref{fig:dictenc}) is composed of a
classifier followed by a dictionary. The classifier splits its input
into $h$ independent \emph{assignments}, each a distribution over one
head's $k$ \emph{entries}, and the dictionary maps each entry to a
representative vector (an \emph{atom}). Each layer writes the
assignment-weighted sum of its atoms to the accumulator; after each
layer the accumulator is decoded and the residual
$\mathbf{x} - \mathrm{dec}(\mathbf{r}_\ell)$ is passed to the next
layer, where $\mathbf{r}_\ell$ accumulates the $\ell$ layers before it.
Deep supervision~\cite{lee2015deeply} constrains every prefix of the
stack to reconstruct the input.

Because entries within a head are mutually exclusive rather than
independent, each head acts as a contrastive set: a $k$-way classifier
whose entry swaps are steering moves. An auxiliary between-head penalty
is intended to push each head onto a disjoint set of dictionary
coordinates (Figure~\ref{fig:head-simplex}), though in practice summed
models do not form a coordinate partition (Section~\ref{sec:disjoint}).

\begin{figure}[tp]
\centering
\resizebox{\linewidth}{!}{\begin{tikzpicture}[
    font=\footnotesize,
    >={Stealth[round,length=1.8mm,width=1.6mm]},
    blk/.style={draw, rounded corners=2pt, align=center,
                minimum height=8mm, inner xsep=4pt},
    hook/.style={draw, rounded corners=2pt, densely dashed, align=center,
                 inner sep=3pt, font=\scriptsize, black!55},
    dec/.style={draw, rounded corners=2pt, fill=black!8,
                minimum height=6.5mm, inner xsep=3pt},
    cond/.style={draw, rounded corners=2pt, fill=black!4, align=center,
                 inner sep=2.5pt, font=\scriptsize},
    add/.style={draw, circle, minimum size=4.5mm, inner sep=0pt},
    lab/.style={inner sep=1.5pt, font=\scriptsize},
    sig/.style={->, semithick},
    aux/.style={->, densely dashed, semithick, black!45},
  ]

  %% ---- layer row --------------------------------------------------
  \node[lab] (x)  at (-1.5,0) {$\mathbf{x}\in\mathbb{R}^{d}$};
  \node[blk] (D0) at ( 0.3,0) {$\mathtt{DictEnc}_0$};
  \node[blk] (D1) at ( 4.3,0) {$\mathtt{DictEnc}_1$};
  \node      (Dd) at ( 7.8,0) {$\cdots$};
  \node[blk] (DL) at ( 9.7,0) {$\mathtt{DictEnc}_{l-1}$};

  \draw[sig] (x) -- (D0);

  \node[hook] (noise) at (-1.5,2.6) {$+\,\varepsilon$};
  \draw[aux] (noise) -- (x);

  %% ---- residual accumulator (top rail, y = 2.6) --------------------
  \node[add] (A0) at ( 2.3,2.6) {$+$};
  \node[add] (A1) at ( 6.1,2.6) {$+$};
  \node[add] (AL) at (11.9,2.6) {$+$};

  \node[lab] (R0) at (0.6,2.6) {$\mathbf{r}_0\!=\!\mathbf{0}$};
  \draw[sig]      (R0) -- (A0);
  \draw[sig]      (A0) -- (A1);
  \draw[sig]      (A1) -- (AL);
  \draw[semithick] (AL) -- (12.35,2.6);

  %\node[lab] at (8.9,3.05) {$\mathbf{r}\in\mathbb{R}^{e}$: residual accumulator};

  %% ---- per-layer contributions g_l f'_l (enter from below) ----------
  \draw[sig] (D0.north) -- ++(0,1.1) -| (A0.south);
  \draw[sig] (D1.north) -- ++(0,1.1) -| (A1.south);
  \draw[sig] (DL.north) -- ++(0,1.1) -| (AL.south);
  \node[lab,anchor=west] at (0.38,1.05) {$g_0\mathbf{f}'_0$};
  \node[lab,anchor=west] at (4.38,1.05) {$g_1\mathbf{f}'_1$};
  \node[lab,anchor=west] at (9.78,1.05) {$g_{l-1}\mathbf{f}'_{l-1}$};

  %% ---- shared decoder, one decode per prefix -----------------------
  \node[dec] (Y1) at ( 2.75,-1.9) {$\mathrm{dec}$};
  \node[dec] (Y2) at ( 6.55,-1.9) {$\mathrm{dec}$};
  \node[dec] (YL) at (12.35,-1.9) {$\mathrm{dec}$};

  \draw[sig] ( 2.75,2.6) -- (Y1.north);
  \draw[sig] ( 6.55,2.6) -- (Y2.north);
  \draw[sig] (12.35,2.6) -- (YL.north);
  \fill ( 2.75,2.6) circle (1.1pt);
  \fill ( 6.55,2.6) circle (1.1pt);

  \node[lab,anchor=west] at ( 2.85,1.0) {$\mathbf{r}_1$};
  \node[lab,anchor=west] at ( 6.65,1.0) {$\mathbf{r}_2$};
  \node[lab,anchor=west] at (12.45,1.0) {$\mathbf{r}_{l}$};

  %% ---- deep supervision --------------------------------------------
  \node[lab] (L1) at ( 0.9,-1.9) {$\mathrm{MSE}(\mathbf{x},\hat{\mathbf{x}}_1)$};
  \node[lab] (L2) at ( 4.9,-1.9) {$\mathrm{MSE}(\mathbf{x},\hat{\mathbf{x}}_2)$};
  \node[lab] (LL) at (10.7,-1.9) {$\mathrm{MSE}(\mathbf{x},\hat{\mathbf{x}}_{l})$};
  \draw[aux] (Y1.west) -- (L1);
  \draw[aux] (Y2.west) -- (L2);
  \draw[aux] (YL.west) -- (LL);

  %% ---- residual forwarding ------------------------------------------
  \node[cond] (C1) at (2.75,-3.9)
  {$\mathrm{sg}\,(\mathbf{x}-\hat{\mathbf{x}}_1)$};
  \node[cond] (C2) at (6.55,-3.9)
  {$\mathrm{sg}\,(\mathbf{x}-\hat{\mathbf{x}}_2)$};
  \draw[sig] (Y1) -- node[right,lab]{$\hat{\mathbf{x}}_1$} (C1);
  \draw[sig] (Y2) -- node[right,lab]{$\hat{\mathbf{x}}_2$} (C2);

  \draw[sig] (C1.east) -- node[left,lab,pos=0.72]{$\mathbf{y}_1$} (D1.south);
  \draw[sig] (C2.east) -- node[left,lab,pos=0.72]{$\mathbf{y}_2$} (7.8,-0.5);

  %% ---- model output --------------------------------------------------
  \node[lab] (out) at (12.35,-3.9) {$\hat{\mathbf{x}}$};
  \draw[sig,thick] (YL) -- node[right,lab]{$\hat{\mathbf{x}}_{l}$} (out);

  %% ---- x broadcast to every conditioning step -------------------------
  \draw[aux] (x.south) -- (-1.5,-5.3) -- (2.75,-5.3) -- (C1.south);
  \draw[aux] (2.75,-5.3) -- (6.55,-5.3) -- (C2.south);
  \fill[black!45] (2.75,-5.3) circle (1.1pt);

\end{tikzpicture}

}
\caption[\texttt{Ontologizer} architecture]{\textbf{\texttt{Ontologizer} architecture.}
	An \texttt{Ontologizer} consists of $l$ stacked \texttt{DictEnc} layers.
	Each returns a pooled feature vector $\mathbf{f}'_\ell\in\mathbb{R}^{e}$,
	which is written, scaled by its input's gain $g_\ell$, to a residual
	accumulator $\mathbf{r}$.
	Layers are indexed $\ell = 0, \dots, l-1$ and accumulator states by the
	number of layers applied: $\mathbf{r}_0 := \mathbf{0}$ and
	$\mathbf{r}_{\ell+1} = \mathbf{r}_\ell + g_\ell\,\mathbf{f}'_\ell$, so
	layer $\ell$ reads state $\ell$ and writes state $\ell+1$; a stop-gradient
	variant ($\mathrm{sg}$ on $\mathbf{r}_\ell$, available but off in every model
	reported here) would prevent gradient communication between layers.
	$\mathrm{dec}$ is a linear decoder with shape $\mathbb{R}^{d \times e}$,
	shared by all layers.
	Deep supervision makes the code a successive refinement, coarse-to-fine
	in distortion as in residual
	quantizers~\cite{juang1982,barnes1996,zeghidour2021} and
	nested-prefix SAEs~\cite{bussmann2025}: the component encoded by
	previous layers is inaccessible to later layers.
	After $\ell$ layers, $\hat{\mathbf{x}}_\ell = \mathrm{dec}(\mathbf{r}_\ell)$, and
	the component left unexplained is
	$\mathbf{y}_\ell = \mathbf{x} - \hat{\mathbf{x}}_\ell$, so $\mathbf{y}_0 = \mathbf{x}$.
	Layer $\ell$ classifies the shape of its input,
	$\mathbf{u}_\ell
	  = \mathbf{y}_\ell / \mathrm{L2}(\mathbf{y}_\ell) \oplus [1]$, and scales
	what it writes by the gain $g_\ell = \mathrm{L2}(\mathbf{y}_\ell)$.
	Normalizing prevents exploding gradients from repeated
	application of bilinear layers, each of which scales its input by
	$\lVert\cdot\rVert^{2}$; the constant coordinate keeps the sign of
	$\mathbf{y}_\ell$, which the bilinear form would otherwise lose.
	$\mathrm{MSE}(\mathbf{x},\hat{\mathbf{x}}_\ell)$ is computed separately for
	each prefix $\ell = 1, \dots, l$. The input-noise hook drawn on $\mathbf{x}$ is disabled in
	every model reported here.
}
\label{fig:onto-stack}
\end{figure}

\subsubsection{Classifier}
\label{sec:bilinear}
Each $\dictenc$'s classifier is a multihead bilinear map
($\bilinearblock$) with weights
$[\mathsf{W}, \mathsf{V}] \in \mathbb{R}^{2 \times h \times k \times d}$
($d+1$ with the constant coordinate of Eq.~\eqref{eq:y_norm}). It
reads the normalized shape $\mathbf{u}_\ell$ of the layer's input (the
residual left by previous layers, or $\mathbf{x}$ at layer 0), written
$\mathbf{x}$ below for a generic input, and returns logits
$K \in \mathbb{R}^{h \times k}$:
\begin{equation}
	\bilinearblock([\mathsf{W},\mathsf{V}],\mathbf{x}) :=
	\left(\sum^{d}_{q=1}W_{\cdot\cdot q}\,x_q\right)
	 \odot \left(\sum^{d}_{q=1}V_{\cdot\cdot q}\,x_q\right).
	 \label{eq:bilinearblock}
\end{equation}
Each logit is a quadratic form in $\mathbf{x}$, the low-rank form of
bilinear classifiers~\cite{pirsiavash2009,kim2017}, so a classifier admits
closed-form eigendecomposition in the manner of bilinear
MLPs~\cite{pearce2025}.

\begin{figure}[tp]
\centering
\resizebox{\linewidth}{!}{\begin{tikzpicture}[
    font=\footnotesize,
    >={Stealth[round,length=1.8mm,width=1.6mm]},
    blk/.style={draw, rounded corners=2pt, align=center,
                minimum height=8mm, inner xsep=4pt},
    hook/.style={draw, rounded corners=2pt, densely dashed, align=center,
                 inner sep=3pt, font=\scriptsize, black!55},
    lab/.style={inner sep=1.5pt, font=\scriptsize, align=center},
    sig/.style={->, semithick, align=center},
    aux/.style={->, densely dashed, semithick, black!45},
  ]

  %% ---- main chain ----------------------------------------------------
  \node[lab] (E)  at (-0.6,0) {$\mathbf{u}$};
  \node[blk] (CL) at ( 1.2,0) {$\bilinearblock^{d \to h \times k}$};
  \node[blk] (SM) at ( 4.9,0) {$\mathrm{softmax}_k(\cdot/t)$};
  \node[blk] (DB) at ( 9.4,0) {$\dictblock^{h \times k \to h \times e}$};
  \node[blk] (SU) at (12.3,0) {$\sum^{h}_{i=1}$};
  \node[lab] (F)  at (13.6,0) {$\mathbf{f}'$\\ pooled\\ features};

  \draw[sig] (E)  -- (CL);
  \draw[sig] (CL) -- node[below,lab]{$K$\\ logits}   (SM);
  \draw[sig] (SM) -- node[below,lab]{$P$\\ classifications}   (DB);
  \draw[sig] (DB) -- node[below,lab]{$F$\\ features} (SU);
  \draw[sig] (SU) -- (F);

  %% ---- training-time hooks and the intervention interface (above) ----
  \node[hook] (nk) at (3.0,1.5) {$+\,\varepsilon_K$,\ winner dropout};
  \node[hook] (iv) at (7.0,2.3) {\texttt{DictIntervention}\\
                                 set / add / sub / ablate / scale};
  \node[hook] (nf) at (11.4,1.5) {$+\,\varepsilon_F$};
  \draw[aux] (nk) -- (3.0,0.20);
  \draw[aux] (iv) -- (7.0,0.20);
  \draw[aux] (nf) -- (11.4,0.20);

  %% ---- where each regulariser attaches (below) -----------------------
  \node[hook] (lp) at (4.9,-1.8)
    {$H$\quad $\mathrm{KL}_m$\quad $\mathrm{cossim}_b$};
  \node[hook] (lf) at (9.4,-1.8)
    {$\mathrm{L1}$\quad $\omega$};
  \draw[aux] (SM.south) -- (lp);
  \draw[aux] (DB.south) -- (lf);

\end{tikzpicture}

}
\caption[\texttt{DictEnc} architecture]{\textbf{\texttt{DictEnc} architecture.}
	Dashed lines above indicate perturbations to the data,
	including noising and causal interventions.
	Dashed lines below indicate statistics used as tuneable auxiliary loss
	functions. Each \texttt{DictEnc} consists of a multihead bilinear
	classifier (Section~\ref{sec:bilinear}) and a multihead dictionary
	(Section~\ref{sec:dict}). Noise $\varepsilon_K$ and $\varepsilon_F$ and winner
	dropout regularize training (Section~\ref{sec:noise}).
	Assignments $P$ are $\mathrm{softmax}_{k}(K/t)$, or its straight-through
	argmax (Section~\ref{sec:labels}); interventions act on $P$. From $P$ we
	collect per-sample entropy $H$, between-sample similarity
	$\mathrm{cossim}_b$, and the divergence of the batch-mean assignment
	from uniform $\mathrm{KL}_m$
	(Eqs.~\eqref{eq:entropy}--\eqref{eq:KL}). A $\dictblock$ maps
	$P \in \mathbb{R}^{h \times k}$ to features
	$F \in \mathbb{R}^{h \times e}$ (Eq.~\eqref{eq:dictblock}); from $P$
	and the atoms we collect the support overlap $\omega$
	(Eq.~\eqref{eq:support}), and from $F$ the sparsity penalty
	$\mathrm{L1}$. $F$ is summed over heads to give pooled
	features $\mathbf{f}' \in \mathbb{R}^{e}$, which are added to the
	accumulator $\mathbf{r}$.
}
\label{fig:dictenc}
\end{figure}

\subsubsection{Assignments and selection}
\label{sec:labels}
Soft assignments are given by $P = \mathrm{softmax}_{k}(K/t)$, with the
temperature $t$ annealed during training to sharpen the distributions.
When trained from scratch with soft assignments, however, the model
continues to rely on the full probability vector rather than its discrete
argmax, since the linear decoder can invert a soft relaxation of
quantization~\cite{agustsson2017,vandenoord2017}
(Section~\ref{sec:prelim-pareto}). To
train a genuinely discrete code, the \emph{straight-through} selection
rule selects the one-hot argmax in the forward pass and uses
$\mathrm{softmax}(K/t)$ as a surrogate gradient in the backward pass,
following the categorical straight-through
estimator~\cite{bengio2013,jang2017}. A hard code transmits
$l\,h\log_2 k$ index bits per sample, with the per-layer gains $g_\ell$
as its only continuous values.

From the assignments $P_n$ of each sample $n$ we collect per-sample
entropy, in bits,
\begin{equation}
	\mathrm{entropy}_k(P)_i :=
	-\sum^{k}_{j=1}p_{ij}\log_2 p_{ij},
	\qquad
	H := \mathbb{E}_{n}\,\mathbb{E}_{i}\!\left[
	\mathrm{entropy}_k(P_n)_i\right],
	\label{eq:entropy}
\end{equation}
between-sample similarity, the mean absolute cosine of the per-head
features over pairs of samples, with
$\cos\angle(\mathbf{a}, \mathbf{b}) :=
\langle\mathbf{a},\mathbf{b}\rangle / (\lVert\mathbf{a}\rVert\lVert\mathbf{b}\rVert)$
the cosine of the angle between two vectors, evaluated through each
head's dictionary Gram $G_i = \lvert D_{i}\rvert \lvert D_{i}\rvert^{\top}$,
\begin{equation}
	\mathrm{cossim}_b :=
	\mathbb{E}_{n,n'}\big\lvert
	\cos\angle(\mathbf{f}_{n}, \mathbf{f}_{n'})
	\big\rvert,
	\qquad
	\langle \mathbf{f}_{n}, \mathbf{f}_{n'}\rangle
	= \sum^{h}_{i=1}\mathbf{p}_{ni}^{\top}G_{i}\,\mathbf{p}_{n'i},
	\label{eq:cossim_b}
\end{equation}
and the divergence of the batch-mean assignment from uniform,
\begin{equation}
	\mathrm{KL}_m := \mathbb{E}_{i}\!\left[\log_2 k -
	\mathrm{entropy}_k\!\left(\mathbb{E}_n[P_n]\right)_i\right].
	\label{eq:KL}
\end{equation}
$\mathrm{KL}_m$ anchors each head's mean assignment to uniform, which
keeps every entry in use.

\subsubsection{Dictionary}
\label{sec:dict}
A $\dictblock$ with weights $\mathsf{D} \in \mathbb{R}^{h \times k \times e}$
maps assignments to per-head features $F \in \mathbb{R}^{h \times e}$,
whose row $\mathbf{f}_i$ is head $i$'s output,
\begin{equation}
	\mathbf{f}_i :=
	\sum^{k}_{j=1} p_{ij}\,\mathbf{a}_{ij},
	\qquad
	\mathbf{a}_{ij} := \lvert \mathsf{D}_{ij\cdot}\rvert .
	\label{eq:dictblock}
\end{equation}
Atoms are non-negative and assignments lie on the simplex, so features
are strictly additive and each head's output lies in the convex hull of
its atoms (Figure~\ref{fig:head-simplex}).

\begin{figure}[tp]
\centering
\begin{tikzpicture}[
    font=\footnotesize,
    >={Stealth[round,length=1.8mm,width=1.6mm]},
    lab/.style={inner sep=1.5pt, font=\scriptsize},
  ]
  %% ---- head i --------------------------------------------------------
  \fill[black!8]   (0,0) -- (2.4,0) -- (1.2,2.0) -- cycle;
  \draw[semithick] (0,0) -- (2.4,0) -- (1.2,2.0) -- cycle;
  \fill (0,0)     circle (1.5pt) node[below,lab]{$\mathbf{a}_{i1}$};
  \fill (2.4,0)   circle (1.5pt) node[below,lab]{$\mathbf{a}_{i2}$};
  \fill (1.2,2.0) circle (1.5pt) node[above,lab]{$\mathbf{a}_{i3}$};
  \fill (1.05,0.72) circle (1.7pt) node[left,lab]{$\mathbf{f}_i$};

  %% ---- head i' ------------------------------------------------------
  \fill[black!8]   (5.0,0) -- (7.4,0) -- (6.2,2.0) -- cycle;
  \draw[semithick] (5.0,0) -- (7.4,0) -- (6.2,2.0) -- cycle;
  \fill (5.0,0)   circle (1.5pt) node[below,lab]{$\mathbf{a}_{i'1}$};
  \fill (7.4,0)   circle (1.5pt) node[below,lab]{$\mathbf{a}_{i'2}$};
  \fill (6.2,2.0) circle (1.5pt) node[above,lab]{$\mathbf{a}_{i'3}$};
  \fill (6.55,0.60) circle (1.7pt) node[right,lab]{$\mathbf{f}_{i'}$};

  \draw[<->, densely dashed, semithick, black!55] (1.35,0.71) -- (6.30,0.61);
  \node[lab] at (3.85,1.35)
    {$\omega$ penalizes heads $i$, $i'$ sharing coordinates};

  \node[lab, anchor=north, align=center] at (3.7,-0.9)
  {$\mathbf{f}'=\sum^h_{i=1} \mathbf{f}_i$,\qquad
  $\mathbf{f}_i=\sum^k_{j=1} p_{ij}\,\mathbf{a}_{ij}$,\qquad $\sum^k_{j=1} p_{ij}=1$};
\end{tikzpicture}


\caption[Per-head geometry]{\texttt{DictBlock} \textbf{head geometry} for $h{=}2$; $k{=}3$.
	The output space of a head $i$ is a simplex with its atoms
	$\mathbf{a}_{ij}$, the rows of $\lvert D_i\rvert$, as vertices. The assignment $P$ is a soft
	choice among mutually contrastive categories rather than a set of
	independently firing features, and because $p_{ij} \in [0,1]$ the
	output lies within the convex hull of the head's atoms.
	Because $\lvert \mathsf{D}\rvert\geq0$ and $P\geq0$, $F\geq0$, so
	$\cos\angle(\mathbf{f}_i, \mathbf{f}_{i'}) \in [0,1]$, and driving the support overlap $\omega$
	to zero encourages the heads to partition the coordinates of $\mathbb{R}^e$
	between them (Section~\ref{sec:disjoint}).
}
\label{fig:head-simplex}
\end{figure}

\paragraph{Disjoint support}
\label{sec:disjoint}
For non-negative vectors the inner product is a sum of non-negative
terms, so it vanishes if and only if the supports are disjoint,
\begin{equation}
	\langle \mathbf{u},\mathbf{v}\rangle = 0
	\iff
	\mathrm{supp}(\mathbf{u}) \cap \mathrm{supp}(\mathbf{v}) = \emptyset .
	\label{eq:disjoint}
\end{equation}
Between heads, orthogonality on non-negative atoms is therefore a
support-partition condition rather than a rotation constraint. The
support overlap $\omega$ measures the mean cosine between heads'
usage-weighted coordinate profiles:
\begin{equation}
	\boldsymbol{\pi}_i := \sum^{k}_{j=1} \bar{p}_{ij}\,
	\lvert \mathbf{a}_{ij}\rvert,
	\qquad
	\omega := \mathbb{E}_{i \neq i'}\,
	\frac{\langle \boldsymbol{\pi}_i, \boldsymbol{\pi}_{i'}\rangle}
	     {\lVert \boldsymbol{\pi}_i\rVert \lVert \boldsymbol{\pi}_{i'}\rVert},
	\qquad
	\bar{\mathbf{p}}_i := \mathbb{E}_{n}\,[\mathbf{p}_{ni}],
	\label{eq:support}
\end{equation}
which vanishes if and only if no two heads share an active coordinate.
Because softmax assignments give every entry positive weight, two heads
that are orthogonal on even a single sample must have disjoint supports
across all $k^2$ pairs of their atoms, making the partition a weight-level
property of $\mathsf{D}$ that holds for all inputs. In the
concatenated-head variant, each head writes to a separate coordinate
slice by construction and $\omega$ is identically zero.

This design distinguishes between-head and within-head geometry: disjoint
support assigns each head its own coordinate block, while the $k$ atoms
\emph{within} a head act as contrastive alternatives spread across that
block. When heads are summed rather than concatenated, $\omega$ is only a
soft penalty, and empirically the within-head contrast emerges while a
between-head coordinate partition does not
(Section~\ref{sec:prelim-orth}).

\subsubsection{Decoder}
\label{sec:decoder}
Each layer's pooled output $\mathbf{f}'_\ell = \sum_{i=1}^{h}\tilde{\mathbf{f}}_{i}$ is scaled
by its input's gain $g_\ell$ and added to the accumulator, which is
projected back to the input space by a single unbiased linear decoder
shared across all layers:
\begin{equation}
	\mathbf{r}_\ell = \sum_{\ell'=0}^{\ell-1} g_{\ell'}\,\mathbf{f}'_{\ell'},
	\qquad
	\hat{\mathbf{x}}_\ell = W_{\mathrm{dec}}\mathbf{r}_\ell,
	\qquad
	W_{\mathrm{dec}} \in \mathbb{R}^{d \times e}.
	\label{eq:dec}
\end{equation}
The residual $\mathbf{y}_\ell$ entering layer $\ell$ is stop-gradiented
so that earlier layers cannot pass a communication code through the
residual instead of reducing reconstruction error. It is decomposed into
a scalar gain $g_\ell$ and a normalized shape vector $\mathbf{u}_\ell$
with a constant coordinate appended, which forms the input to layer
$\ell$'s classifier:
\begin{equation}
	\mathbf{y}_\ell = \mathrm{sg}\!\left[\mathbf{x} - \hat{\mathbf{x}}_\ell\right],
	\qquad
	\mathbf{u}_\ell =
	\frac{\mathbf{y}_\ell}{\mathrm{L2}(\mathbf{y}_\ell) + \epsilon} \oplus [1],
	\qquad
	g_{\ell} := \mathrm{sg}\!\left[\mathrm{L2}(\mathbf{y}_\ell)\right].
	\label{eq:y_norm}
\end{equation}
At layer 0 the empty sum gives $\mathbf{r}_0 = \mathbf{0}$ and
$\mathbf{y}_0 = \mathbf{x}$, so the first layer classifies the input
directly. When gain scaling (\texttt{resid\_gain}) is enabled, each
layer classifies the direction of the current residual and scales its
write by $g_\ell$; otherwise $g_\ell = 1$.

\paragraph{Direct atoms and fibers.}
We also evaluate two architectural variants of the dictionary. In the
\emph{direct-atom} variant, the shared decoder is omitted and signed
atoms live directly in the $d$-dimensional output space, so each atom is
an explicit vector in activation space. In the \emph{fiber} variant, each
entry carries a local rank-$r$ basis and each head reads $r$ continuous
coordinates from its input, augmenting the discrete entry choice with a
low-rank continuous correction; training fibers on top of a frozen hard
code refines reconstruction without altering the discrete assignments.

\subsubsection{Noise and dropout}
\label{sec:noise}
Classifier logits take Gaussian noise $\varepsilon_K$, scaled relative to
the logit standard deviation under hard selection; dictionary features
take Gaussian noise $\varepsilon_F$ proportional to each coordinate's
batch standard deviation, with the scale stop-gradiented so the model
cannot suppress noise by shrinking feature variance. Winner dropout sets
a head's argmax logit to $-\infty$ with probability $p_{\mathrm{drop}}$.

\subsection{Training}
\label{sec:training}

\subsubsection{Embedding model}
\label{sec:sonar}
We embed texts using SONAR~\cite{sonar}, a multimodal and
language-agnostic sentence embedding model consisting of a
translation-trained encoder (mean-pooled into a 1024-dimensional vector)
and a paired decoder that generates text from that vector in a chosen
target language. Because translations map to nearby vectors, reconstructed
or steered embeddings can be decoded into a fixed language and re-encoded
to measure text fidelity and cycle-consistent steering. Large Concept
Models~\cite{lcmteam2024largeconceptmodelslanguage} generate sentence
embeddings directly in SONAR space, making SONAR a natural testbed for
embedding-level interventions, though SONAR decoding is known to be
sensitive to small perturbations~\cite{lcmteam2024largeconceptmodelslanguage}.
As a second substrate where sparse dictionary learning is well
established and reconstructions can be spliced directly into a forward
pass, we also train models on the residual stream of GPT-2
small~\cite{radford2019}.

\subsubsection{Loss function}
\label{sec:loss}
Reconstruction is trained with a whitened mean-squared error, where
$\mathbf{w}$ is the per-dimension inverse variance normalized to mean one,
\begin{equation}
	\mathrm{MSE}(X,\hat{X}) := \frac{1}{bd}\sum^{b}_{n=1}\sum^{d}_{q=1}
	w_q\left(\hat{x}_{nq} - x_{nq}\right)^2,
	\label{eq:mse}
\end{equation}
applied with deep supervision~\cite{lee2015deeply} at every prefix
$\hat{X}_\ell$ ($\ell = 1, \dots, l$) and averaged across layers, so each
layer is trained to explain the residual left by earlier layers. The
auxiliary statistics are weighted by $\mathbf{s}$ and summed across
layers:
\begin{equation}
	L := \frac{1}{l}\sum^{l}_{\ell=1}\mathrm{MSE}(X,\hat{X}_\ell)
	+ \sum^{l-1}_{\ell=0}\mathbf{s}\cdot\boldsymbol{\psi}_\ell,
	\qquad
	\boldsymbol{\psi}_\ell := \big[\, \mathrm{L1} \;\; \mathrm{cossim}_b \;\;
	\mathrm{KL}_m \;\; \omega \,\big]^{\top}_\ell,
	\label{eq:loss_total}
\end{equation}
where $\mathrm{L1}$ is the batch mean of the feature L1 norm. Because
minimizing $\mathrm{cossim}_b$ also reduces per-sample entropy $H$, we
set the weight on $H$ to zero. Most models reported below were trained
with the predecessor of $\omega$ (the within-sample mean cosine between
heads' feature vectors); Section~\ref{sec:prelim-orth} reports a direct
sweep over the $\omega$ penalty on GPT-2.

\subsection{Evaluation}
\label{sec:eval}
We evaluate Ontologizer models against $\topk$ SAEs~\cite{gao2024}
trained on the same cached embeddings under the same whitened MSE
objective, reporting whitened fraction of variance unexplained
($\mathrm{FVU}_w$) on a held-out evaluation tail. A structural-ablation
ladder of SAE variants introduces Ontologizer components one step at a
time; two rungs that add per-head competition over a shared decoder are
equivalent to single-layer Ontologizers with linear classifiers
(\texttt{g160top1} with hard selection and \texttt{g160softmax} with
softmax) and are treated as simplified model variants rather than
external baselines. Our evaluation suite measures:
\begin{itemize}
	\item reconstruction error against code capacity and a Gaussian
	rate--distortion baseline~\cite{cover2006};
	\item within-head semantic coherence of decoded atoms against
	size-matched random subsets from the same dictionary;
	\item steering efficacy and collateral in embedding space (how often
	a step along a target direction realizes the target entry and how
	many other heads it flips, relative to a random step of equal norm),
	along with additivity of two-head interventions;
	\item cross-seed reproducibility of head partitions and output
	contributions; and
	\item blinded autointerp detection.
\end{itemize}
Because AxBench~\cite{wu2025axbenchsteeringllmssimple} found that
prompting, fine-tuning, and supervised difference-of-means vectors
outperform SAE latents at steering while linear probes do not, we compare
our unsupervised steering primitives directly against both
difference-of-means and linear-probe vectors.
